\scanpage{17}
\mspage{A}
\[
\begin{aligned}
(a)&\quad -\ 1\to N\to G\to H\to1
 &&\text{ext. de groupes top.}\quad [\,N\text{ discret}\,]\\
(b)&\quad \deleted{\unsure{\varphi}}\quad 1\to\pi\to\widetilde H\to H\to1
 &&\text{ext. de groupes top.}\quad [\,\pi\text{ \rlap{\(\big/\)}discret}\,]\\
(c)&\quad 1\to\widetilde H_0\to\widetilde H\to\unsure{\mathfrak S|\mathfrak G}\to1
 &&\text{ext. de gr. top.}\quad [\,\mathfrak S\text{ discret}\,]
\end{aligned}
\]

On suppose donné (b) et (c) \inserted{et $N$}, on s'intéresse aux ext. $G$ de $H$ par $N$, \uline{munies} d'une trivialisation \inserted{\uline{stricte} (\unsure{isom.|triv.|i.e. comm.} avec pdt. pr. direct)} de l'image inverse de cette extension par $\widetilde H_0\to H$ (composé $\widetilde H_0\hookrightarrow\widetilde H\to H$) i.e. d'un hom $\widetilde H_0\to G$ qui remonte $\widetilde H_0\to H$ et centralise $N$

(\uline{NB} $N$ étant discret, si $\widetilde H_0$ est connexe et simplement connexe, il existe pour $G$ donné une unique telle trivialisation\ldots)

\inserted{On suppose de plus que ce splitting est invariant par les opérations de $\widetilde G$ sur $\widetilde H_0$ (via $\widetilde H$ \unsure{qui} $\deleted{\unsure{\to H}}$) et sur $G$ (via $\widetilde G\to G$) (cf ci-dessous pour déf. de $\widetilde G$)}
Soit $\widetilde G$ l'image inverse de $G$ par $\widetilde H\to H$
\[
1\to N\to\widetilde G\to\widetilde H\to1
\]
La donnée d'une ext $G$ est équivalente à celle de $\widetilde G$, \uline{munie} d'une trivialisation \inserted{stricte} sur $\pi$. La \deleted{trivialisation de} structure supplémentaire sur l'ext. $G$ (trivialisation de l'image inverse par $\widetilde H_0\to H$) équivaut alors à une trivialisation \inserted{stricte} de $\widetilde G$ sur le sous-groupe $\widetilde H_0$. Donc au résumé on doit étudier les ext. $\widetilde G$ de $\widetilde H$ par $N$, trivialisées \inserted{strictemt} simultanément sur $\pi$ \uline{et} sur $\widetilde H_0$ [rappelons nous que cette dernière structure, \deleted{\unsure{triv.}} \unsure{fixant|spécifique|sur} $\widetilde H_0$, est automatique dans les conditions dites]. Or une ext. \deleted{\unsure{trivial.}} \inserted{$\widetilde G$ de} $\widetilde H$ par $N$ trivialisée \inserted{\unsure{strit|strict|strictement}} sur $\widetilde H_0$ équivaut à une ext. de $\mathfrak S$ par $N$, soit $\widetilde G_0$ :
\[
1\to N\to\overset{\deleted{\unsure{\sim}}}{G}_0\to\mathfrak S\to1
\]
($\widetilde G$ est l'image inverse par $\widetilde H\to\mathfrak S$). De plus, il faut donner un splittage \inserted{strict} du diagramme
% archive sheet number: 16
